\section{Comparación de los tres procedimientos}
\label{sec:comparacion}

\subsection{MoSCoW y votación ponderada}

MoSCoW clasificó veinte requisitos en nueve Must, siete Should y cuatro Could; ninguno recibió categoría Won't. Los valores comunicados por el equipo corresponden a 44 puntos de esfuerzo Must, 32 Should y 16 Could, con un total de 92 puntos relativos. El grupo Must representa el 47,83\,\% y cumple la restricción de la ecuación~\eqref{eq:must}.

Se asignaron rangos medios a los requisitos empatados dentro de cada grupo MoSCoW: Must ocupa las posiciones 1 a 9 (rango 5), Should las posiciones 10 a 16 (rango 13) y Could las posiciones 17 a 20 (rango 18,5). La votación se ordenó por puntos descendentes, también con rangos medios cuando las puntuaciones coinciden. Con los veinte pares de rangos, la correlación de Spearman, calculada como correlación de Pearson entre rangos, fue \(\rho=0{,}6605\). Se trata de una asociación positiva apreciable, pero no de una equivalencia entre las técnicas. No se infiere causalidad ni se generaliza a una población escolar: la votación recoge tres participantes.

\begingroup
\small
\begin{longtable}{@{}p{0.14\linewidth}p{0.21\linewidth}rrr@{}}
\caption{Contraste entre MoSCoW y votación}\label{tab:comparacion-moscow-votos}\\
\toprule
\textbf{ID} & \textbf{MoSCoW} & \textbf{R. MoSCoW} & \textbf{Puntos} & \textbf{R. votos}\\
\midrule
\endfirsthead
\toprule
\textbf{ID} & \textbf{MoSCoW} & \textbf{R. MoSCoW} & \textbf{Puntos} & \textbf{R. votos}\\
\midrule
\endhead
\bottomrule
\endlastfoot
RF-04 & Must & 5,0 & 31 & 1,0\\
RF-08 & Must & 5,0 & 25 & 2,0\\
RF-05 & Must & 5,0 & 24 & 3,0\\
RF-01 & Must & 5,0 & 21 & 4,0\\
RF-10 & Must & 5,0 & 19 & 5,0\\
RF-07 & Must & 5,0 & 18 & 6,5\\
RF-12 & Could & 18,5 & 18 & 6,5\\
RF-09 & Should & 13,0 & 17 & 8,0\\
RF-02 & Must & 5,0 & 16 & 9,0\\
RNF-03 & Must & 5,0 & 15 & 10,0\\
RF-06 & Should & 13,0 & 13 & 11,0\\
RF-11 & Should & 13,0 & 12 & 12,5\\
RNF-02 & Should & 13,0 & 12 & 12,5\\
RF-03 & Must & 5,0 & 11 & 14,5\\
RNF-01 & Should & 13,0 & 11 & 14,5\\
RNF-04 & Could & 18,5 & 9 & 16,5\\
RNF-05 & Could & 18,5 & 9 & 16,5\\
RNF-06 & Should & 13,0 & 8 & 18,0\\
RNF-07 & Could & 18,5 & 6 & 19,0\\
RNF-08 & Should & 13,0 & 5 & 20,0\\
\end{longtable}
\endgroup


Entre las coincidencias, RF-04 (matrícula, 31 puntos), RF-08 (registro de calificaciones, 25) y RF-05 (control de cupos, 24) son Must y ocupan las tres primeras posiciones de votación. En cambio, RF-12 (reporte de matrícula) recibe 18 puntos y comparte el rango 6,5 en la votación, aunque está clasificado como Could. RF-03 es Must, pero obtiene 11 puntos y comparte el rango 14,5. Estas diferencias justifican una negociación del alcance guiada por dependencias y esfuerzo, no solamente por la suma de puntos.

\subsection{Kano frente a la votación}

Kano devuelve una categoría para ocho características, no un ranking de veinte requisitos. Se observaron dos características indiferentes (K01 y K08), cuatro unidimensionales (K02 a K05) y dos atractivas (K06 y K07). En la tabla siguiente se muestran solo esas ocho correspondencias, sin imputar una categoría a los doce requisitos no evaluados por Kano.

\begingroup
\small
\begin{longtable}{@{}p{0.14\linewidth}p{0.40\linewidth}rrc@{}}
\caption{Características Kano frente a la votación}\label{tab:comparacion-kano-votos}\\
\toprule
\textbf{ID} & \textbf{Característica} & \textbf{Puntos} & \textbf{Rango} & \textbf{Kano}\\
\midrule
\endfirsthead
\toprule
\textbf{ID} & \textbf{Característica} & \textbf{Puntos} & \textbf{Rango} & \textbf{Kano}\\
\midrule
\endhead
\bottomrule
\endlastfoot
RF-04 & Matrícula & 31 & 1,0 & I\\
RF-06 & Comprobante de matrícula & 13 & 11,0 & O\\
RF-08 & Registro de calificación & 25 & 2,0 & O\\
RF-10 & Consulta de notas & 19 & 5,0 & O\\
RF-11 & Boletín académico & 12 & 12,5 & O\\
RNF-01 & Rendimiento de matrícula & 11 & 14,5 & A\\
RNF-04 & Facilidad de uso de matrícula & 9 & 16,5 & A\\
RNF-05 & Compatibilidad móvil & 9 & 16,5 & I\\
\end{longtable}
\endgroup


RF-04 fue el requisito mejor puntuado de la votación (31 puntos), pero su característica K01 obtuvo tres de cinco respuestas indiferentes. RF-08 alcanzó 25 puntos y su característica K03 fue unidimensional. Esta diferencia metodológica no se interpreta como contradicción de las personas participantes, porque las técnicas preguntan por aspectos diferentes. El tamaño de las muestras (cinco formularios Kano y tres votantes) restringe la interpretación a los registros analizados. En la tabulación Kano publicada se repitió el código P04 y no aparece P03; la copia exportada proporcionada por el equipo sí registra P03. No se alteran los archivos de respuestas de otros integrantes en esta sección.
